%--------------------------------------------------------------------------------------
% Feladatkiiras (a tanszeken atveheto, kinyomtatott valtozat)
%--------------------------------------------------------------------------------------
\clearpage
\begin{center}
\large
\textbf{Ökoszisztéma szimulációja Unity3D alapokon}\\
\end{center}

A szimulációk lehetővé teszik a valós rendszerek viselkedésének biztonságos és költséghatékony vizsgálatát. Segítségükkel komplex folyamatokat modellezhetünk, különböző változókkal kísérletezhetünk, és megérthetjük a valós rendszerek működését anélkül, hogy a valóságban beavatkoznánk. Az ökoszisztéma-szimulációk különösen fontosak, mivel bemutatják a természetes egyensúly, a táplálkozási láncok és a populációdinamika összefüggéseit vizuálisan ábrázolható grafikonok formájában is. A Unity keretrendszer jó alapot biztosít egy ilyen szimuláció elkészítéséhez.

A hallgató feladata a korábban általa kidolgozott, Unity3D alapú ökoszisztéma szimulációs rendszer továbbfejlesztése, hogy az minél kiterjeszthetőbb legyen és faj-specifikus tulajdonságok és viselkedések legyenek benne. A részletes statisztikák vizuális megjelenítése szintén fontos szempont. Összefoglalva, a hallgató feladatának a következőkre kell kiterjednie:

\begin{itemize}
    \item Röviden mutassa be a Unity keretrendszert és a korábban elkészült szimulációs megoldást!
    \item Fejlessze tovább a megoldást, hogy az minél kiterjeszthetőbb legyen! Ügyeljen a fontos események logolására is!
    \item Vezessen be faj-specifikus tulajdonságokat és viselkedéseket és vizsgálja meg, ezek hogyan befolyásolják a szimulációt!
    \item Készítsen részletes statisztikákat a szimulációk futása után! Gondoskodjon ezek vizuális megjelenítéséről is!
\end{itemize}